\chapter{Hỏi Để Giải Tỏa Căng Thẳng}
\chapterintro{Không phải câu hỏi nào cũng để lấy kiến thức}{Có những câu hỏi được sinh ra chỉ để lớp thở nhẹ hơn và quay lại học.}

\section{Nếu mô tả tuần này bằng một từ, em chọn từ nào?}
\begin{cauhoibox}{Câu hỏi}
Nếu mô tả tuần này bằng một từ, em chọn từ nào?
\end{cauhoibox}
\begin{taisaocoich}
Câu hỏi một từ giúp học sinh nói ra trạng thái mà không phải kể lể dài. Rất tốt khi lớp đang mệt hoặc đang chất nhiều áp lực.
\end{taisaocoich}
\begin{goiybien}
Cho 5 đến 7 em nói nhanh, không phân tích quá sâu nếu đang trong tiết học chính.
\end{goiybien}
\ngancach

\section{Điều gì làm em thấy đỡ nhất trong một ngày học nhiều việc?}
\begin{cauhoibox}{Câu hỏi}
Điều gì làm em thấy đỡ nhất trong một ngày học nhiều việc?
\end{cauhoibox}
\begin{taisaocoich}
Câu này mở ra hướng tích cực, không dồn lớp vào than thở mà kéo các em về những điểm tựa nhỏ nhưng thật.
\end{taisaocoich}
\begin{goiybien}
Rất hợp cho đầu tuần căng, sau kiểm tra, hoặc giờ chủ nhiệm.
\end{goiybien}
\ngancach

\section{Nếu được bớt một thứ cho lớp mình hôm nay, em bớt gì?}
\begin{cauhoibox}{Câu hỏi}
Nếu được bớt một thứ cho lớp mình hôm nay, em bớt gì?
\end{cauhoibox}
\begin{taisaocoich}
Câu hỏi này vừa vui, vừa giúp lớp nói ra thứ đang gây mệt mà không biến thành than phiền quá mức.
\end{taisaocoich}
\begin{goiybien}
Có thể chốt bằng câu hỏi ngược: \textit{``Vậy mình thêm gì để lớp dễ thở hơn?''}
\end{goiybien}
\ngancach

\section{Nếu lớp mình cần một nút thở, nó nằm ở đâu?}
\begin{cauhoibox}{Câu hỏi}
Nếu lớp mình cần một nút thở ngay lúc này, nó nên nằm ở đâu?
\end{cauhoibox}
\begin{taisaocoich}
Câu hỏi có hình ảnh nên học sinh trả lời nhanh và thật. Nó cũng giúp thầy cô đo xem lớp đang muốn giảm áp ở khâu nào: nội dung, nhịp, hay không khí.
\end{taisaocoich}
\begin{goiybien}
Có thể gợi ý 3 lựa chọn: ở câu hỏi, ở tốc độ, hay ở cách làm việc. Sau đó chọn một điểm để chỉnh.
\end{goiybien}
\ngancach

\section{Điều gì lớp mình nên buông xuống trong 30 giây tới?}
\begin{cauhoibox}{Câu hỏi}
Điều gì lớp mình nên buông xuống trong 30 giây tới để dễ học tiếp hơn?
\end{cauhoibox}
\begin{taisaocoich}
Thay vì chỉ hỏi lớp đang mệt gì, câu này kéo học sinh sang hành động nhẹ và gần. Nó rất hợp để cắt căng giữa tiết.
\end{taisaocoich}
\begin{goiybien}
Câu trả lời có thể là: bỏ bớt nói chuyện riêng, bỏ bớt sợ sai, bỏ bớt nhìn điện thoại trong đầu. Nghe xong là chuyển tiết ngay.
\end{goiybien}
